\documentclass[10pt]{scrartcl}
\usepackage[mdthm,colorsec,noauthor,noasy]{evan}
\usepackage{handout}
\newcommand{\ip}[2]{\langle #1,#2\rangle}
\DeclareMathOperator{\tr}{tr}\DeclareMathOperator{\nul}{null}\DeclareMathOperator{\range}{range}\DeclareMathOperator{\spn}{span}
\newcommand{\LL}{\mathcal L}\newcommand{\PP}{\mathcal P}\newcommand{\MM}{\mathcal M}
\title{Math 110 --- Week 10}
\subtitle{Orthogonality (Notes Ch.~18, LADR 6B--6C, HW10)}
\author{}\date{}\ihead{}
\begin{document}
\maketitle
\section{Facts used from earlier weeks}
$V$ is finite-dimensional over $\FF$ ($\RR$ or $\CC$), $n=\dim V$. Shorthand: ``$S$ kills $v$'' means $Sv=0$, and ``$q$ kills $T$'' means $q(T)=0$. From Week 9: inner products and norms, the orthogonal decomposition $u=cv+w$ with $c=\ip uv/\|v\|^2$ and $w\perp v$ (LADR 6.13), Pythagoras, and Cauchy--Schwarz.

\begin{enumerate}[label=(T\arabic*),leftmargin=2.6em,itemsep=8pt]
  \ii \tool{Rank--nullity.} For linear $T\colon V\to W$: $\dim V=\dim\nul T+\dim\range T$.
  \why{\begin{steps}
    \ii Extend a basis $u_1,\dots,u_k$ of $\nul T$ to a basis $u_1,\dots,u_k,v_1,\dots,v_r$ of $V$. It suffices to show $Tv_1,\dots,Tv_r$ is a basis of $\range T$.
    \ii Spanning: $T$ of any vector is a combination of the $Tu_i$ (which are $0$) and the $Tv_j$.
    \ii Independent: if $\sum c_jTv_j=0$, then $\sum c_jv_j\in\nul T$ is a combination of the $u_i$, and independence of the basis of $V$ gives all $c_j=0$.
  \end{steps}}
  \ii \tool{Dimension shortcuts.}
  \begin{enumerate}[(a)]
    \ii An independent list of $n$ vectors is a basis.
    \ii A subspace of dimension $n$ is all of $V$.
    \ii If $U\cap W=\{0\}$ and $\dim U+\dim W=n$, then $V=U\oplus W$.
    \ii $\dim(U+W)=\dim U+\dim W-\dim(U\cap W)$.
  \end{enumerate}
  \why{\begin{steps}
    \ii (a), (b): an independent list extends to a basis, and every basis has length $n$.
    \ii (c): a basis of $U$ followed by a basis of $W$ is independent (a relation gives $u=-w\in U\cap W=\{0\}$) and has length $n$, so it is a basis by (a).
    \ii (d): extend a basis $x_1,\dots,x_k$ of $U\cap W$ to a basis $x_i,u_j$ of $U$ and, separately, to a basis $x_i,w_l$ of $W$. Together the $x_i,u_j,w_l$ span $U+W$; it remains to show they are independent.
    \ii If $x+u+w=0$ ($x,u,w$ combinations of the $x_i$, $u_j$, $w_l$), then $w=-x-u\in U\cap W$, so $w$ is a combination of the $x_i$. The basis $x_i,w_l$ of $W$ is independent, so $w=0$. Then $x+u=0$, and independence of $x_i,u_j$ gives $x=u=0$.
  \end{steps}}
  \ii \tool{Products of operators.} If $S,T$ are injective, so is $ST$. So if a product sends some $v\ne0$ to $0$, some factor is not injective.
  \why{$STv=0\Rightarrow Tv=0\Rightarrow v=0$.}
  \ii \tool{Polynomials in $T$ commute.} $p(T)q(T)=q(T)p(T)$. Hence $\nul p(T)$ and $\range p(T)$ are $T$-invariant.
  \why{Both products equal $(pq)(T)$. If $p(T)v=0$, then $p(T)(Tv)=T\,p(T)v=0$. And $T\bigl(p(T)u\bigr)=p(T)(Tu)\in\range p(T)$.}
  \ii \tool{Fundamental theorem of algebra.} A nonconstant $p\in\PP(\CC)$ factors as $c(z-\lambda_1)\cdots(z-\lambda_m)$ (cited). Over any field, $p(\lambda)=0\iff(z-\lambda)\mid p$.
  \why{Divide $p$ by $z-\lambda$: the remainder is a constant, and evaluating at $\lambda$ shows it equals $p(\lambda)$.}
  \ii \tool{Eigenvalues.}
  \begin{enumerate}[(a)]
    \ii $\lambda$ is an eigenvalue $\iff T-\lambda I$ is not injective.
    \ii If $Tv=\lambda v$, then $p(T)v=p(\lambda)v$.
    \ii Eigenvectors for distinct eigenvalues are independent (Prop.~54).
    \ii Over $\CC$ with $V\ne0$, every operator has an eigenvalue.
  \end{enumerate}
  \why{\begin{steps}
    \ii (a) is the definition: $Tv=\lambda v$ with $v\ne0$ means $(T-\lambda I)v=0$ with $v\ne0$.
    \ii (b): $T^kv=\lambda^kv$ by induction; add up the terms of $p$.
    \ii (d): take $v\ne0$. The $n+1$ vectors $v,Tv,\dots,T^nv$ are dependent, so $p(T)v=0$ for some $p\ne0$; $p$ is not constant since $v\ne0$.
    \ii Factor $p=c(z-\lambda_1)\cdots(z-\lambda_m)$ (T5). The product $c(T-\lambda_1I)\cdots(T-\lambda_mI)$ sends $v\ne0$ to $0$, so some $T-\lambda_jI$ is not injective (T3).
  \end{steps}}
  \ii \tool{Upper-triangular matrices.} $\MM(T)$ is upper-triangular in $v_1,\dots,v_n$ $\iff$ every $\spn(v_1,\dots,v_j)$ is invariant. Then the eigenvalues of $T$ are the diagonal entries. Over $\CC$ every operator has an upper-triangular matrix.
  \why{\begin{steps}
    \ii Column $j$ holds the coordinates of $Tv_j$, so ``upper-triangular'' means $Tv_j\in\spn(v_1,\dots,v_j)$ for every $j$.
    \ii If $a_{jj}=\lambda$: $T-\lambda I$ maps $\spn(v_1,\dots,v_j)$ into $\spn(v_1,\dots,v_{j-1})$, a smaller space, so it is not injective (T1). So $\lambda$ is an eigenvalue.
    \ii If no $a_{jj}$ equals $\lambda$: $\MM(T-\lambda I)$ is triangular with nonzero diagonal. Its columns are independent: in a relation, the right-most column used is the only one with a nonzero entry in its own diagonal row. So $T-\lambda I$ is injective.
    \ii Existence over $\CC$, by induction on $n$: take an eigenvalue $\lambda$ (T6(d)). $U=\range(T-\lambda I)$ is invariant (T4) and $\dim U<n$, since $T-\lambda I$ is not injective (T1).
    \ii By induction $U$ has a basis $u_1,\dots,u_m$ with every $\spn(u_1,\dots,u_j)$ invariant. Extend it to a basis $u_1,\dots,u_m,v_1,\dots,v_k$ of $V$.
    \ii $Tv_j=(T-\lambda I)v_j+\lambda v_j\in U+\spn(v_j)\subseteq\spn(u_1,\dots,u_m,v_1,\dots,v_j)$. So every span of an initial part of the list is invariant.
  \end{steps}}
\end{enumerate}

\section{This week}
\subsection{Orthogonality (Ch.~18)}

\begin{qbox}
  In a general basis, finding coordinates means solving a linear system. If the basis vectors are perpendicular unit vectors, each coordinate is one inner product. We want such bases, a way to build them, and to use them to split spaces into perpendicular pieces.
\end{qbox}

\begin{definition*}[Def.~46]
  $u\perp v$ if $\ip uv=0$. $e_1,\dots,e_m$ is \vocab{orthonormal} if $\ip{e_i}{e_j}=\delta_{ij}$.
\end{definition*}

\begin{proposition*}[Prop.~84--86 / LADR 6.24, 6.25, 6.30]
  For an orthonormal list: $\|\sum a_ie_i\|^2=\sum|a_i|^2$, and the list is independent. For an orthonormal basis:
  \[ v=\sum_i\ip v{e_i}e_i,\qquad\|v\|^2=\sum_i|\ip v{e_i}|^2. \]
\end{proposition*}
\begin{ideabox}
Norm: Pythagoras handles two perpendicular vectors, and $a_me_m$ is perpendicular to $a_1e_1+\dots+a_{m-1}e_{m-1}$. So add one term at a time.

Independence: If $\sum a_ie_i=0$, its norm $\sqrt{\sum|a_i|^2}$ is $0$, so every $a_i=0$.

Coordinates: To find $a_k$ in $v=\sum a_ie_i$, take the inner product with $e_k$: the terms with $i\ne k$ give $0$, and the $k$-th gives $a_k$. This is the reason to want orthonormal bases: coordinates without solving a linear system.
\end{ideabox}
\begin{proof}
  \begin{steps}
    \ii $e_m\perp a_1e_1+\dots+a_{m-1}e_{m-1}$, so Pythagoras adds $|a_m|^2$; induct.
    \ii $\sum a_ie_i=0\Rightarrow\sum|a_i|^2=0\Rightarrow$ all $a_i=0$.
    \ii If $v=\sum a_ie_i$, then $\ip v{e_k}=\sum a_i\ip{e_i}{e_k}=a_k$.
  \end{steps}
\end{proof}
\needspace{7\baselineskip}\textbf{Parseval} (infinite-dimensional version): for $f(x)=\sum c_ne^{2\pi inx}$ on $[0,1]$, $\sum|c_n|^2=\int_0^1|f|^2$. With $f(x)=x=\frac12-\frac1{2\pi i}\sum_{n\ne0}\frac1ne^{2\pi inx}$ this gives $\frac14+\frac1{4\pi^2}\sum_{n\ne0}\frac1{n^2}=\frac13$, so $\sum_{n\ge1}\frac1{n^2}=\frac{\pi^2}6$.

\begin{theorem*}[Prop.~87 / LADR 6.32: Gram--Schmidt]
  For independent $v_1,\dots,v_n$: $e_1=v_1/\|v_1\|$, and for $j\ge2$
  \[ w_j=v_j-\sum_{k<j}\ip{v_j}{e_k}e_k,\qquad e_j=w_j/\|w_j\|. \]
  Then $e_1,\dots,e_n$ is orthonormal and $\spn(e_1,\dots,e_j)=\spn(v_1,\dots,v_j)$ for every $j$.
\end{theorem*}
\begin{ideabox}
To show: Orthonormal $e_1,\dots,e_n$ with $\spn(e_1,\dots,e_j)=\spn(v_1,\dots,v_j)$ for every $j$.

For a unit vector $e$, $v-\ip ve\,e\perp e$ (LADR 6.13 with $\|e\|=1$). Subtracting a combination of $e_1,\dots,e_{j-1}$ from $v_j$ does not change the span of $e_1,\dots,e_{j-1},v_j$.

Build one vector at a time: from $v_j$ subtract its component $\ip{v_j}{e_k}e_k$ along each earlier $e_k$. Since the earlier $e_k$ are perpendicular to each other, the subtraction for $e_k$ does not change the inner products with the other $e_i$. Then divide by the length.

Only risk: $w_j=0$, which would put $v_j$ in $\spn(v_1,\dots,v_{j-1})$; independence rules it out.
\end{ideabox}
\begin{proof}
  \begin{steps}
    \ii If $w_j=0$, then $v_j\in\spn(e_1,\dots,e_{j-1})=\spn(v_1,\dots,v_{j-1})$, contradicting independence.
    \ii For $k<j$: $\ip{w_j}{e_k}=\ip{v_j}{e_k}-\ip{v_j}{e_k}=0$. And $\|e_j\|=1$.
    \ii $e_j\in\spn(v_1,\dots,v_j)$, so the spans are nested; both have dimension $j$, so they are equal.
  \end{steps}
\end{proof}
\begin{itemize}
  \ii Every finite-dimensional inner product space has an orthonormal basis; every orthonormal list extends to one (extend to a basis, then Gram--Schmidt; the first vectors do not change) (LADR 6.35, 6.36).
  \ii An orthonormal list of length $\dim V$ is an orthonormal basis (Prop.~85, T2).
  \ii If the $v_j$ are dependent, $w_j=0$ exactly when $v_j\in\spn(v_1,\dots,v_{j-1})$: skip it and continue.
  \ii \textbf{Schur (LADR 6.38):} over $\CC$, $T$ is upper-triangular in some \emph{orthonormal} basis. Take a triangular basis (T7) and apply Gram--Schmidt: it keeps the spans $\spn(v_1,\dots,v_j)$, which are invariant, so the matrix stays triangular.
\end{itemize}

\begin{example*}[Ex.~70: Gram--Schmidt on $\PP_2$, $\ip pq=\int_0^1p\bar q$]
  \begin{steps}
    \ii $\|1\|=1$: $e_1=1$.
    \ii $w_2=x-\ip x1=x-\frac12$, $\|w_2\|^2=\frac1{12}$: $e_2=\sqrt3(2x-1)$.
    \ii $w_3=x^2-\frac13-\frac12(2x-1)=x^2-x+\frac16$, $\|w_3\|=\frac1{6\sqrt5}$: $e_3=\sqrt5(6x^2-6x+1)$.
  \end{steps}
  Order of work: $w_j$, then $\|w_j\|^2$, then divide. Check $w_j\perp$ the earlier $e$'s. (These are the shifted Legendre polynomials; on $[-1,1]$, HW10 P1, the Legendre polynomials.)
\end{example*}

\begin{proposition*}[Prop.~88 / LADR 6.42, 6.58: Riesz representation]
  For finite-dimensional $V$, every $\psi\in V'$ is $\psi(u)=\ip uv$ for a unique $v$. In an orthonormal basis, $v=\sum_j\overline{\psi(e_j)}e_j$.
\end{proposition*}
\begin{ideabox}
To show: A vector $v$ with $\psi(u)=\ip uv$ for all $u$.

On $\RR^n$ a functional is $\psi(u)=a\cdot u$ for a row vector $a$, so $v=a$. A linear functional is determined by its values on a basis.

Write the unknown $v=\sum a_je_j$ in an orthonormal basis and test on $e_i$: $\ip{e_i}v=\bar a_i$ must equal $\psi(e_i)$. So $a_i=\overline{\psi(e_i)}$, and this $v$ works because both sides agree on a basis.
\end{ideabox}
\begin{proof}
  \begin{steps}
    \ii With $v=\sum\overline{\psi(e_j)}e_j$: $\ip{e_i}v=\psi(e_i)$. Two linear functionals equal on a basis are equal.
    \ii Unique: if $\ip uv=\ip u{v'}$ for all $u$, take $u=v-v'$ to get $\|v-v'\|^2=0$.
  \end{steps}
  (The notes use $\Gamma(v)=\ip-v$: real-linear, not complex-linear since $\Gamma(iv)=-i\Gamma(v)$, injective, and $\dim_\RR V=\dim_\RR V'$.)
\end{proof}

\begin{definition*}[Def.~47 / LADR 6.46, 6.47]
  $U^\perp=\{v:\ip vu=0\text{ for all }u\in U\}$, a subspace (linearity in the first slot).
\end{definition*}
\begin{itemize}
  \ii $\{0\}^\perp=V$; $V^\perp=\{0\}$; $U\cap U^\perp=\{0\}$ (a vector perpendicular to itself is $0$).
  \ii $U\subseteq W\Rightarrow W^\perp\subseteq U^\perp$.
  \ii $(U+W)^\perp=U^\perp\cap W^\perp$ (perpendicular to all $u+w$ iff to all $u$ and all $w$).
  \ii $(U\cap W)^\perp=U^\perp+W^\perp$ (finite dim., HW10 P4): take $\perp$ of $(U^\perp+W^\perp)^\perp=U\cap W$ (Prop.~90).
  \ii Finite dim.: $U^\perp=\{0\}\iff U=V$; $U^\perp=V\iff U=\{0\}$ (Prop.~89).
  \ii \textbf{Computing $U^\perp$:} require $\ip v{u_i}=0$ for a spanning list of $U$ and solve the linear system. Under Riesz, $U^\perp$ matches the annihilator $U^0$.
\end{itemize}

\begin{theorem*}[Prop.~89 / LADR 6.49, 6.50]
  If $U$ is finite-dimensional, $V=U\oplus U^\perp$. If $V$ is finite-dimensional, $\dim U^\perp=\dim V-\dim U$.
\end{theorem*}
\begin{ideabox}
To show: Every $v$ is $u+w$ with $u\in U$ and $w\perp U$, in exactly one way.

For a line $U=\spn(e)$, $\|e\|=1$, this is LADR 6.13: $u=\ip ve\,e$.

For a subspace, take an orthonormal basis $e_1,\dots,e_m$ of $U$ and use all the components: $u=\sum\ip v{e_k}e_k$. Then $\ip{v-u}{e_k}=0$ for every $k$, so $v-u\perp U$. Uniqueness: $U\cap U^\perp=0$, since a vector perpendicular to itself is $0$.

Dimension: $\dim U^\perp=\dim V-\dim U$ then follows from the direct sum.
\end{ideabox}
\begin{proof}
  \begin{steps}
    \ii Direct: if $u\in U\cap U^\perp$, then $\ip uu=0$, so $u=0$.
    \ii Take an orthonormal basis $e_1,\dots,e_m$ of $U$ (Gram--Schmidt). Set $u=\sum\ip v{e_k}e_k\in U$, $w=v-u$.
    \ii $\ip w{e_k}=\ip v{e_k}-\ip v{e_k}=0$ for each $k$, so $w\perp U$, and $v=u+w$.
  \end{steps}
\end{proof}

\begin{proposition*}[Prop.~90 / LADR 6.52]
  If $U$ is finite-dimensional, $(U^\perp)^\perp=U$.
\end{proposition*}
\begin{ideabox}
To show: $(U^\perp)^\perp=U$. The inclusion $U\subseteq(U^\perp)^\perp$ is immediate from the definitions; the work is showing that $(U^\perp)^\perp$ is not bigger.

Finite dimension: Count: $\dim(U^\perp)^\perp=n-(n-\dim U)=\dim U$ (Prop.~89 twice), and a subspace of the same dimension is equal (T2).

Take $v\in(U^\perp)^\perp$ and split $v=u+w$ (Prop.~89). Then $w\in U^\perp$, and also $w=v-u\in(U^\perp)^\perp$. So $w$ is perpendicular to itself, and $w=0$.
\end{ideabox}
\begin{proof}
  $U\subseteq(U^\perp)^\perp$: every $u\in U$ is perpendicular to every element of $U^\perp$. Conversely let $v\in(U^\perp)^\perp$, $v=u+w$ ($u\in U$, $w\in U^\perp$). Then $w=v-u\in(U^\perp)^\perp$ and $w\in U^\perp$, so $\ip ww=0$ and $w=0$.
\end{proof}

\begin{definition*}[Def.~48 / LADR 6.55]
  For finite-dimensional $U$, the \vocab{orthogonal projection} is $P_U(u+w)=u$ ($u\in U$, $w\in U^\perp$). It is the unique projector with $\range P=U$, $\nul P=U^\perp$.
\end{definition*}
\begin{itemize}
  \ii \textbf{Uniqueness:} if $P^2=P$, $\range P=U$, $\nul P=U^\perp$, then $u=Px$ gives $Pu=P^2x=u$, so $P(u+w)=u$.
  \ii \textbf{Properties} (LADR 6.57), all from $v=P_Uv+(v-P_Uv)$:
  \begin{itemize}
    \ii $P_U$ is linear; $P_U^2=P_U$; $\range P_U=U$; $\nul P_U=U^\perp$; $v-P_Uv\in U^\perp$.
    \ii $\|P_Uv\|\le\|v\|$ (Pythagoras).
    \ii $I-P_U=P_{U^\perp}$ if $V$ is finite-dimensional ($V=U^\perp\oplus U$).
  \end{itemize}
  \ii \textbf{Any projector} ($P^2=P$) has $V=\nul P\oplus\range P$ (Prop.~91). Idea: $Pv$ is the obvious range part, and $P(v-Pv)=Pv-P^2v=0$. If $w=Pv\in\nul P$, then $w=P^2v=Pw=0$.
  \ii It is orthogonal only when $\nul P\perp\range P$: $P(x,y)=(0,y)$ is $P_{\spn(e_2)}$; $P(x,y)=(x+y,0)$ is not orthogonal.
\end{itemize}

\begin{proposition*}[Prop.~92 / LADR 6.57(i)]
  If $e_1,\dots,e_m$ is an orthonormal basis of $U$: $P_Uv=\sum_k\ip v{e_k}e_k$.
\end{proposition*}
\begin{ideabox}
To show: The coordinates of $P_Uv$ in an orthonormal basis $e_1,\dots,e_m$ of $U$.

Coordinates in an orthonormal basis are inner products: $P_Uv=\sum\ip{P_Uv}{e_k}e_k$ (Prop.~86 in $U$).

$v-P_Uv\perp e_k$, so $\ip v{e_k}=\ip{P_Uv}{e_k}$. The coordinates can be computed from $v$, without knowing $P_Uv$ first.
\end{ideabox}
\begin{proof}
  $v=u+w$: $\ip v{e_k}=\ip u{e_k}$, so $P_Uv=u=\sum\ip u{e_k}e_k=\sum\ip v{e_k}e_k$ (Prop.~86 in $U$).
\end{proof}

\begin{proposition*}[HW10 P5, LADR 6.61: best approximation]
  For $u\in U$: $\|v-P_Uv\|\le\|v-u\|$, equality only at $u=P_Uv$. So the distance from $v$ to $U$ is $\|v-P_Uv\|$.
\end{proposition*}
\begin{ideabox}
To show: Among $u\in U$, $\|v-u\|$ is smallest at $u=P_Uv$.

Split the error: $v-u=(v-P_Uv)+(P_Uv-u)$. The first part is in $U^\perp$ and does not depend on $u$; the second is in $U$. They are perpendicular, so $\|v-u\|^2=\|v-P_Uv\|^2+\|P_Uv-u\|^2$, which is smallest exactly when the second term is $0$.
\end{ideabox}
\begin{proof}
  $v-u=(v-P_Uv)+(P_Uv-u)$ with pieces in $U^\perp$ and $U$, so $\|v-u\|^2=\|v-P_Uv\|^2+\|P_Uv-u\|^2$.
\end{proof}
\begin{fig}
\begin{tikzpicture}[scale=1.15,>=Stealth]
  \draw[thick] (-0.6,0) -- (4.6,0) node[right] {$U$};
  \fill (2,1.8) circle (1.8pt) node[above] {$v$};
  \fill (2,0) circle (1.8pt) node[below] {$P_Uv$};
  \fill (3.8,0) circle (1.8pt) node[below] {$u$};
  \draw[OliveGreen,thick] (2,0) -- (2,1.8);
  \node[OliveGreen!70!black,left] at (2,0.9) {$v-P_Uv\perp U$};
  \draw[red,thick,dashed] (3.8,0) -- (2,1.8);
  \node[red,right] at (2.95,1.0) {$v-u$};
  \draw (2,0.18) -- (2.18,0.18) -- (2.18,0);
\end{tikzpicture}
\figcap{$v-u$ is the hypotenuse of a right triangle. One leg, $v-P_Uv$, does not depend on $u$; the other, $P_Uv-u$, lies in $U$ and is $0$ only when $u=P_Uv$.}
\end{fig}

\begin{itemize}
  \ii \textbf{Bessel:} for any orthonormal $e_1,\dots,e_m$, $\sum|\ip v{e_k}|^2\le\|v\|^2$. The left side is $\|P_Uv\|^2$ with $U=\spn(e_k)$.
  \ii \textbf{HW10 P7:} a projector is orthogonal iff $\|Pv\|\le\|v\|$ for all $v$. Idea for $\Leftarrow$: suppose $u\in\range P$ and $w\in\nul P$ are not perpendicular. The point of the line $u+tw$ closest to $0$ has $t=-\ip uw/\|w\|^2\ne0$ and is shorter than $u$. But $P(u+tw)=u$, so $\|P(u+tw)\|>\|u+tw\|$. (Details in the solutions.)
  \ii \textbf{Computing $P_Uv$:} (1) Gram--Schmidt on $U$, then Prop.~92. (2) Or write $P_Uv=\sum c_iu_i$ for a basis $u_i$ and solve $\ip{v-\sum c_iu_i}{u_j}=0$ for all $j$. Check that $v-P_Uv\perp$ each $u_j$.
\end{itemize}

\begin{example*}[Ex.~71]
  $P_U(x^3)$, $U=\spn(1,x)$, $\ip pq=\int_0^1pq$. With $e_1=1$, $e_2=\sqrt3(2x-1)$: $\ip{x^3}{e_1}=\frac14$, $\ip{x^3}{e_2}=\frac{3\sqrt3}{20}$, so $P_U(x^3)=\frac14+\frac9{20}(2x-1)=\frac9{10}x-\frac15$. Check: $x^3-\frac9{10}x+\frac15$ is perpendicular to $1$ and $x$.
\end{example*}

\begin{trapbox}
  \ii Orthogonal vectors are independent only if nonzero.
  \ii $P^2=P$ does not make $P$ orthogonal: $P(x,y)=(x+y,0)$ has $\|P(1,1)\|=2>\sqrt2$.
  \ii $V=U\oplus U^\perp$ can fail for infinite-dimensional $U$: in $\ell^2$, the finitely supported sequences have $U^\perp=\{0\}$. Also $\ell^2$ has no orthonormal (Hamel) basis.
  \ii Gram--Schmidt depends on the order and on the inner product; normalize only after subtracting.
  \ii Riesz over $\CC$: the coefficients are conjugated.
\end{trapbox}

\clearpage

\section{Problem set}

\begin{prob}
\textbf{HW10 P1. }\ On $\PP_2(\RR)$ use the inner product $\ip pq=\int_{-1}^1p(x)q(x)\,dx$. Apply the Gram--Schmidt procedure to the ordered basis $1,x,x^2$ to obtain an orthonormal basis of $\PP_2(\RR)$.
\end{prob}

\begin{prob}
\textbf{HW10 P2. }\ Let $W=\spn\big((1,-1,1,0),(0,2,1,0)\big)\subseteq\RR^4$ with the standard inner product. Find a basis of $W^\perp$.
\end{prob}

\begin{prob}
\textbf{HW10 P3. }\ On $\PP_2(\RR)$ use the same inner product as in Problem 1, and let $W=\spn(x+1)$. Find a basis of $W^\perp$.
\end{prob}

\begin{prob}
\textbf{HW10 P4. }\ Let $U$ and $W$ be subspaces of a finite-dimensional inner-product space $V$. Prove $(U\cap W)^\perp=U^\perp+W^\perp$.
\end{prob}

\begin{prob}
\textbf{HW10 P5. }\ Let $U$ be a subspace of a finite-dimensional inner-product space $V$, and let $e_1,\dots,e_m$ be an orthonormal basis of $U$.
\begin{enumerate}[(a)]
  \ii Prove that $P_Uv=\sum_{j=1}^m\ip v{e_j}e_j$.
  \ii Prove the best-approximation property $\|v-P_Uv\|\le\|v-u\|$ for every $u\in U$, and determine when equality holds.
\end{enumerate}
\end{prob}

\begin{prob}
\textbf{HW10 P6. (optional)}\ On $\PP_1(\RR)$ use the inner product $\ip fg=\int_{-1}^1f(x)g(x)\,dx$. Find the unique polynomial $p\in\PP_1(\RR)$ such that $\int_0^2f(x)\,dx=\ip fp$ for every $f\in\PP_1(\RR)$.
\end{prob}

\begin{prob}
\textbf{HW10 P7. (optional)}\ Let $P$ be a projector on a finite-dimensional inner-product space, so $P^2=P$. Prove that $P$ is an orthogonal projection if and only if $\|Pv\|\le\|v\|$ for every $v$.
\end{prob}

\begin{prob}
\textbf{O1. }\ Find the polynomial $p\in\PP_1(\RR)$ that minimizes $\int_0^1\big(x^2-p(x)\big)^2dx$, and the minimum value.
\end{prob}

\begin{prob}
\textbf{O2. }\ In $\RR^4$ with the dot product let $U=\spn\big((1,1,1,1)\big)$ and $v=(1,2,3,4)$. Find $P_Uv$, a basis of $U^\perp$, and the distance from $v$ to $U$.
\end{prob}

\begin{prob}
\textbf{O3. Reflections.}\ Let $U$ be a subspace of a finite-dimensional inner product space and $R=I-2P_U$. Prove that $\|Rv\|=\|v\|$ for all $v$ and that $R^2=I$.
\end{prob}

\end{document}
